\section{Optimalizace}\label{sec:optimalizace}

Po přední části máme ověřenou reprezentaci programu. Mohli bychom z~ní rovnou vytvořit cílový kód,~ale výsledek by často prováděl zbytečnou práci.

\emph{Optimalizace programu} je úprava jeho vnitřní reprezentace,~která zachová pozorovatelné chování programu a~zlepší zvolené vlastnosti výsledného kódu,~například dobu běhu,~velikost nebo spotřebu energie.

Slovo \uv{optimalizace} neznamená,~že překladač vždy nalezne nejlepší možný program. Takový požadavek by byl příliš drahý a~obecně i~algoritmicky problematický. Překladač používá sadu transformací,~které bývají výhodné a~jejichž správnost umí zdůvodnit.

\subsection{Lokální optimalizace}

\emph{Základní blok} je souvislá posloupnost instrukcí,~do níž se vstupuje pouze na začátku a~z~níž se odchází pouze na konci. Uvnitř bloku se tedy tok řízení nevětví. \emph{Lokální optimalizace} pracuje s~jedním takovým blokem.

Mezi běžné lokální úpravy patří:
\begin{itemize}
    \item \emph{výpočet konstant} (anglicky \emph{constant folding}) --- výraz $2\cdot60$ lze spočítat už při překladu;
    \item \emph{šíření konstant a~kopií} --- známe-li $x=120$,~můžeme následující použití $x$ nahradit touto hodnotou;
    \item \emph{odstranění mrtvého kódu} --- výpočet,~jehož výsledek se nikde nepoužije,~může být zbytečný;
    \item \emph{odstranění společných podvýrazů} --- nezměnily-li se operandy,~nemusíme stejný výraz počítat podruhé.
\end{itemize}

\begin{figure}[h]
    \centering
    \begin{tikzpicture}[
    code/.style={task,minimum width=46mm,minimum height=23mm,align=left,font=\ttfamily\small}
]
    \node[code,fill=red!8] (before) at (0,0)
        {x := 2 * 60\\y := x + 5\\print y};
    \node[code,fill=green!10] (after) at (7.2,0)
        {x := 120\\y := 125\\print 125};
    \draw[diredge] (before)--node[above,align=center,font=\small]
        {výpočet konstant\\a jejich šíření}(after);
\end{tikzpicture}

    \caption{Příklad výpočtu a~šíření konstant.}
    \label{fig:optimalizace-konstant}
\end{figure}

\warningbox{Transformace musí respektovat přesnou sémantiku jazyka. Nahrazení $0\cdot E$ nulou například není bezpečné,~pokud vyhodnocení $E$ může mít vedlejší účinek nebo skončit chybou.}

\subsection{Výpočet konstant v~AST}

Optimalizátor v~programu \ref{lst:optimizer} prochází strom zdola nahoru. Jsou-li oba operandy binární operace číselné konstanty,~nahradí celý podstrom jediným číselným vrcholem. Dělení nulou záměrně ponechá beze změny,~aby příslušná chyba vznikla až při vykonání programu.

\lstinputlisting[
    caption={Optimalizace výpočtem konstant.},
    label={lst:optimizer}
]{04-kompilatory/code/mini_compiler/optimizer.py}

Ve výrazu $2+3\cdot4$ se nejprve podstrom $3\cdot4$ nahradí číslem $12$ a~potom celý součet číslem $14$. Optimalizace tedy přirozeně využívá stromovou strukturu vytvořenou parserem.

\subsection{Globální optimalizace a~graf toku řízení}

Optimalizace přes podmínky,~cykly a~více bloků potřebuje znát možné pořadí vykonávání programu. Základní bloky proto použijeme jako vrcholy \emph{grafu toku řízení}. Hrana $B_i\to B_j$ znamená,~že po bloku $B_i$ může následovat blok $B_j$.

\begin{figure}[h]
    \centering
    \begin{tikzpicture}
    \node[task,fill=blue!10,minimum width=28mm] (b1) at (0,0) {$B_1$\\podmínka};
    \node[task,fill=orange!12,minimum width=28mm] (b2) at (-2.6,-2) {$B_2$\\větev ano};
    \node[task,fill=orange!12,minimum width=28mm] (b3) at (2.6,-2) {$B_3$\\větev ne};
    \node[task,fill=green!10,minimum width=28mm] (b4) at (0,-4) {$B_4$\\pokračování};
    \node[task,fill=gray!10,minimum width=28mm] (dead) at (5.8,-4) {$B_5$\\mrtvý blok};

    \draw[diredge] (b1)--node[left,font=\small]{ano}(b2);
    \draw[diredge] (b1)--node[right,font=\small]{ne}(b3);
    \draw[diredge] (b2)--(b4);
    \draw[diredge] (b3)--(b4);
\end{tikzpicture}

    \caption{Příklad grafu toku řízení s~nedosažitelným blokem $B_5$.}
    \label{fig:graf-toku-rizeni}
\end{figure}

Blok je dosažitelný,~jestliže do něj vede cesta z~počátečního bloku. Nedosažitelné bloky nelze za běhu navštívit a~můžeme je odstranit. K~jejich nalezení stačí BFS ze sekce \smartref{\showrefs}{sec:bfs}{o~BFS} nebo DFS ze sekce \smartref{\showrefs}{sec:dfs}{o~DFS}.

Další důležitou informací je \emph{živost proměnné}. Hodnota proměnné je za danou instrukcí živá,~pokud může být později přečtena dříve,~než ji přepíše jiná hodnota. Instrukce,~která zapisuje mrtvou hodnotu a~nemá jiný účinek,~může být odstraněna.

Ve zjednodušeném základním bloku lze živost spočítat průchodem instrukcí odzadu:
\begin{enumerate}
    \item začneme množinou proměnných potřebných na konci bloku;
    \item proměnnou,~do které aktuální instrukce zapisuje,~z~množiny odstraníme;
    \item proměnné,~které instrukce čte,~do množiny přidáme.
\end{enumerate}
U~celého grafu se informace mezi bloky předávají a~výpočet se opakuje,~dokud se množiny nepřestanou měnit.

\subsection{Generování bajtkódu}

Po optimalizaci převedeme AST do jednoduchého zásobníkového bajtkódu. Instrukce \texttt{PUSH} vloží konstantu na zásobník,~\texttt{LOAD} načte proměnnou a~instrukce \texttt{ADD},~\texttt{SUB},~\texttt{MUL} a~\texttt{DIV} odeberou dva operandy a~vloží výsledek. Příkazy završují instrukce \texttt{STORE} a~\texttt{PRINT}.

\begin{figure}[h]
    \centering
    \begin{tikzpicture}[node distance=6mm]
    \node[task,fill=blue!10,minimum width=24mm,minimum height=12mm,
        font=\normalsize] (source) {zdrojový\\kód};
    \node[task,fill=orange!12,minimum width=24mm,minimum height=12mm,
        font=\normalsize,right=of source] (front) {přední část\\překladače};
    \node[task,fill=orange!12,minimum width=24mm,minimum height=12mm,
        font=\normalsize,right=of front] (ast) {strom po\\optimalizaci};
    \node[task,fill=green!10,minimum width=24mm,minimum height=12mm,
        font=\normalsize,right=of ast] (bytecode) {bajtkód};
    \node[task,fill=green!10,minimum width=24mm,minimum height=12mm,
        font=\normalsize,right=of bytecode] (vm) {virtuální\\stroj};
    \draw[diredge] (source)--(front);
    \draw[diredge] (front)--(ast);
    \draw[diredge] (ast)--(bytecode);
    \draw[diredge] (bytecode)--(vm);
\end{tikzpicture}

    \caption{Překlad malého jazyka do bajtkódu pro virtuální stroj.}
    \label{fig:preklad-bajtkod}
\end{figure}

Generátor v~programu \ref{lst:codegen} prochází výraz v~pořadí \emph{levý podstrom,~pravý podstrom,~operace}. Tím automaticky vytvoří pořadí vhodné pro zásobníkový stroj.

\lstinputlisting[
    caption={Generování jednoduchého zásobníkového bajtkódu.},
    label={lst:codegen}
]{04-kompilatory/code/mini_compiler/codegen.py}

Po optimalizaci se příkaz \texttt{let x = 2 + 3 * 4;} přeloží jen na
\[
    \texttt{PUSH 14},\qquad \texttt{STORE x}.
\]
Bez optimalizace by vznikly tři instrukce \texttt{PUSH},~násobení,~sčítání a~teprve potom uložení.
